\section{Introduction and Experimental Objective}
\label{sec:intro}

Even at rest, the interval between two consecutive heartbeats changes from one beat to the next. In healthy people these changes are correlated across time scales, so that the pattern of fluctuation over a few beats resembles the pattern over minutes and hours \citep{peng1995,goldberger2002}. This fractal organisation is thought to reflect the interplay of many regulatory loops acting on the sinoatrial node at once, and it degrades in characteristic ways with disease and ageing \citep{lipsitz1992,goldberger2002}. The beat-to-beat timing of the heart therefore offers a continuous, non-invasive readout of how the nervous system is regulating the body.

Heart rate variability is the measure of choice here for practical and theoretical reasons. It is recorded at the resolution of single heartbeats, with a chest strap that tolerates movement, so it can follow a person through a movement practice that would be impossible inside a scanner or with a scalp EEG cap. It reflects the output of a central autonomic network whose cortical nodes overlap with those of attention and self-regulation \citep{benarroch1993,thayer2000}, which makes it a bodily trace of central regulation. And it is already measured by consumer wearables, so any reliable relation between a practice and HRV can be carried out of the laboratory into everyday use.

Most research on heart rate variability (HRV) quantifies how much the heartbeat varies, using indices such as the root mean square of successive differences (RMSSD). These amplitude indices are informative at rest but collapse towards a floor as soon as physical activity rises, losing resolution exactly where a movement task places the participant \citep{gronwald2020biomarker}. The short-term scaling exponent of detrended fluctuation analysis (\aone) quantifies how this variation is organised over windows of 4 to 16 beats. \al{} sits near 1.0 in rested, healthy adults, falls towards 0.75 around the first ventilatory threshold and towards 0.5 at high exercise intensity \citep{rogers2021threshold,rogers2022update}. Because it keeps a dynamic range under load and can be estimated from about two minutes of data \citep{chen2002,rogers2022update}, it is a candidate marker for short bouts of physically light but attentionally demanding movement.

One such class of movement is skilled object manipulation. Juggling, poi, staff and stick work are widely practised, and practitioners describe them as absorbing, calming and restorative, often in the vocabulary of flow \citep{csikszentmihalyi1990}. In cognitive neuroscience, object manipulation already serves as a model of skill learning: learning to juggle changes grey and white matter within weeks to months \citep{draganski2004,scholz2009}. The state of the body during manipulation itself has received far less attention.

These practices belong to a community usually called the \emph{flow arts}, in which the manipulation of props such as poi, staffs, hoops and sticks is practised as a movement discipline, a performance art and, for many, a contemplative practice. The present work grows out of seks.design, an experience-design practice that combines flow-arts movement with \emph{somatic journaling}: the structured recording of bodily sensations alongside reports of mental state. Its guiding idea, called here \emph{somatic interfacing}, is that signals from the body can serve as an interface for tracking, and eventually shaping, mental states. The longer-term aim is a model of conscious experience grounded in embodiment, in which movement, bodily signals and reported experience are mapped onto one another \confirm{wording of the seks.design description}. This thesis takes one step in that direction by asking whether a specific flow-arts movement leaves a measurable trace in cardiac regulation. If a practice reliably induces a distinct state of engagement, that state should leave a trace in autonomic regulation during the practice itself, and \al{} offers a way to look for it.

A second motivation comes from the study of rhythmic coordination. During walking, running and cycling, breathing tends to lock into stable frequency ratios with the movement cycle, and an external rhythm such as a metronome can stabilise this locomotor--respiratory coupling and lower its energetic cost \citep{hoffmann2012,bardy2015}. Music and rhythmic sound, in turn, modulate arousal and can entrain cardiovascular rhythms \citep{bernardi2009,trappe2010}. A rhythmic upper-limb movement therefore offers a plausible route from movement to breathing to the heart, and one in which mechanical coupling plays a smaller role than in locomotion. Whether such coordination behaves like a skill, weak in novices and stable in experts, is the broader question that gives this thesis its title.

The present study uses an isolation stick, a prop of about 50~cm with a weight at each end that produces a rhythmic figure-eight motion. Isolation sticks were created more than 20 years ago by the French artist and manipulation specialist Alexandre ``Pich'' of Le Cirque Graphique, drawing on the work of Michael Moschen on rhythm and motion in object manipulation \todo{source for the history of the isolation stick}. Participants performed two families of movement with it. \emph{Isolations} keep one point of the stick fixed in space while the hand moves around it; they are slow, precise and visually guided. \emph{Spins} rotate the stick continuously around the hand; they are faster and more rhythmic. RR intervals were recorded with a Polar H10 chest strap and streamed to a laptop, where an open pipeline computed \al{} alongside heart rate and amplitude indices of HRV for each two-minute trial and for a two-minute standing baseline in which participants held the stick still.

The experimental objective was to establish whether two minutes of stick manipulation shift \al{} relative to standing still with the stick, and in which direction. Because \al{} falls with exercise intensity, a shift could simply reflect physical load. A second objective was therefore to separate these two accounts by using mean heart rate as a covariate: an effect of manipulation counts as specific only if it exceeds what heart rate alone predicts. A third objective was to test whether the shift depends on the movement performed and on the participant's expertise with the stick, since the same visible movement can correspond to effortful, explicitly controlled action in a novice and to largely automatic action in an experienced practitioner \citep{fitts1967}.

Testing and validating the measurement pipeline itself constitutes a further objective. The capture system was built for this study from consumer hardware and open software, and the real-time \al{} reported by the streaming application can be compared directly with the estimate computed offline. Establishing how far these agree, and how robust \al{} is to the movement artifacts that object manipulation introduces, determines whether the stick can serve as a portable tool for studying states of skilled engagement outside the laboratory.

The thesis is organised as follows. Section~\ref{sec:background} reviews cardiac autonomic control, the method and interpretation of \aone, the evidence for acute changes in \al, rhythmic coordination and its links to the heart, the psychophysiology of skill and flow, and the validity of chest-strap recording, and closes with the hypotheses. Section~\ref{sec:methods} describes participants, design and the capture system; Section~\ref{sec:analysis} the offline signal processing and statistics. Sections~\ref{sec:results} to \ref{sec:conclusion} report, discuss and summarise the results.
